\documentclass[11pt]{article}
\usepackage{amsmath,amssymb,amsthm}
\usepackage[margin=1in]{geometry}

\newcommand{\BRC}{\mathcal{BRC}}
\newcommand{\CE}{\mathcal{E}}
\newcommand{\CR}{\mathcal{R}}
\newcommand{\sch}{\mathfrak{S}}

\begin{document}

\section*{Phantom terms, quantum annihilating pairs, and positivity}

\texttt{full\_CEM} for a permutation $u$ produces a dictionary
$\{R \mapsto \text{cem\_dict}\}$ where each $R$ is an RC graph for~$u$,
and $\text{cem\_dict}$ is its Grassmannian factorization in~$\CE$.
However, some tensors in the CEM dict, when squashed via~$\CR$, land on
an RC graph for a \emph{different} permutation $u' \neq u$.  These are
the \textbf{phantoms}---they live in
$\ker\!\bigl(\CR - \text{(expected perm)}\bigr)$.

When two CEM elements are multiplied, a phantom tensor can pair with a
legitimate tensor and produce:
\begin{itemize}
  \item A \textbf{bad} RC graph (wrong permutation or wrong weight) that
        gets killed, and simultaneously
  \item A \textbf{good} RC graph (correct permutation, correct weight)
        that fills a gap.
\end{itemize}
These come in pairs---a \textbf{quantum annihilating pair}---where the
phantom contribution simultaneously subtracts a term that should not be
present and adds one that should.

This is why the crystal filter works despite the apparent signs: the
phantoms are not creating cancellation in the coefficient~$c_{uv}^w$.
Instead, they perform a \textbf{crystal-theoretic correction}, swapping
contributions that landed in the wrong crystal component for ones in the
right component.  The net effect on~$c_{uv}^w$ is always non-negative
because:
\begin{enumerate}
  \item Without phantoms, the total count would be correct but
        distributed incorrectly across crystal components.
  \item The phantoms redistribute: $-1$ from a forbidden component,
        $+1$ to a missing one.
  \item After filtering by \texttt{is\_highest\_weight} and
        \texttt{extremal\_weight == trimcode}, each surviving component
        has been corrected to its true (non-negative) multiplicity.
\end{enumerate}

A proof would therefore need to show that phantom pairs always act as
\textbf{transpositions on the set of crystal components}: each phantom
involvement swaps exactly one bad contribution for one good one,
preserving the total and maintaining non-negativity after the crystal
filter.

Formally, the phantoms should generate an action on the
fiber~$\CR^{-1}(w)$ that is compatible with the crystal structure, and
the quantum annihilating pairs are precisely the orbits of this action
that cross crystal component boundaries.

\end{document}
